% Nouveau EGA V (cote 25): archive cote 25, site page 22 (PDF page 23).
% typescript. Draft compared: none.
% AI-transcribed by Claude Opus 5.5 (claude-opus-5-5, adaptive thinking, effort high), Message Batches API, batch ega25b (msgbatch_01A7n5YuXubUPNK1zsepi8RG), 2026-11-13; not yet checked by eye.
\scanpage{25/22}
\mspage{\unsure{\hand{19}|19}}

5. \uline{Th\'eor\`emes du type Seidenberg}

5.1.
\restoredfrom{25/21}{https://grothendieck.umontpellier.fr/web/viewer.html?file=../25.pdf\#page=22}{%
Dans \illegible{1} \hand{le} présent n$^o$, nous donnons des conditions \unsure{moyennant} les
quelles l'ensemble E envisagé dans 4.1. est ouvert.  Il
s'agira ici de propriétés P de nature locale sur X resp. $Y_{\hand{\xi}}$ , de sorte
qu'on peut définir l'ensemble U des $y \hand{\in} \underline{Y}$ tels que (si \hand{$\xi$} désigne l'image
de y dans $\hand{\check{}}P$ ) $Y_{\hand{\xi}}$ satisfasse à P en le point y (resp. $G_{\hand{\xi}}$ satisfasse à
P en y). Nous donnerons d'abord des critères pour que U soit ouvert,
en utilisant le par.12.  Comme en tous cas on a
\[ E = P - h(\underline{Y} - U) , \]
il s'ensuit que si U est ouvert et X propre sur S (donc h propre et
à fortiori fermé car $\underline{Y} \to \unsure{Y}$ propre sur X et $\overset{\hand{\vee}}{P}$ séparé sur S), alors E est également ouvert.
}

5.2. Comme nous avons vu au n\textsuperscript{o}1, \underline{Y} est d\'efini dans $X x_S \overset{\hand{\vee}}{P} = X_{\overset{\hand{\vee}}{P}}$
comme le ``diviseur'' d'une section $\phi$ de $\underline{O}_X(1)\otimes_S \underline{O}_{\overset{\hand{\vee}}{P}}(1)$, qui induit pour
tout $\hand{\xi} \hand{\in} \overset{\hand{\vee}}{P}$ une sedtion $\phi_{\hand{\xi}}$ de  $\underline{O}_X(1)\otimes_{k(s)}\underline{O}_{\overset{\hand{\vee}}{P}}(1)(\hand{\xi})$
(faisceau d'ailleurs isomorphe non canoniquement \`a $\underline{O}_X(1)\otimes_{k(s)}k(\hand{\xi}) =$
$\underline{O}_{X_{k(\hand{\xi})}}(1)$ ), de sorte que $Y_{\hand{\xi}}$ n'est autre que le  ``diviseur'' de cet
te section (NB le diviseur d'une section $\phi$ d'un module inversible \underline{L} \placedmargin{$/\!/$}
est d\'efini comme le sous-pr\'esch\'ema ferm\'e d\'efini par l'Id\'eal image de
$\phi^{-1} : \underline{L}^{-1} \hand{\longrightarrow} \underline{O}$ ). Si \underline{F} est un Module sur X, alors son image inverse \underline{G}
sur \underline{Y} , i.e. l'image inverse de $\underline{F}\otimes_{\underline{O}_S}\underline{O}_{\overset{\hand{\vee}}{P}} = \underline{F}_{\overset{\hand{\vee}}{P}}$ sur le sous-pr\'esch\'ema \underline{Y} de
$X_{\overset{\hand{\vee}}{P}}$ , n'est autre que le conoyau du homomorphisme
\[ \phi^{-1}\otimes \mathrm{id}_{\underline{F}_{\overset{\hand{\vee}}{P}}} : \underline{F}_{\overset{\hand{\vee}}{P}}(-1,-1) \mathrel{\hand{-}} F_{\overset{\hand{\vee}}{P}} , \]
o\`u la notation (-1,-1) s'explique d'elle-m\^eme, comme dirait Mike. De
m\^eme, $G_{\hand{\xi}}$ est le conayau de l'homomorphisme analogue
\[ F_{k(\hand{\xi})}(-1,-1) \hand{\longrightarrow} F_{k(\hand{\xi})} \]
o\`u $\hand{\xi}$ est un point de $\overset{\hand{\vee}}{P}$ (et bien entendu on a une interpr\'etation corres-%
pondante lorsque $\hand{\xi}$, au lieu d'\^etre un point de $\overset{\hand{\vee}}{P}$ , d\'esigne un point
de $\overset{\hand{\vee}}{P}$ \`a valeurs dans un S' sur S . ). \hand{\unsure{(}}De fa\c{c}on g\'en\'erale, si L est un
Module inversible quelque part, $\phi$ une section, d\'efinissant le sous-pr\'esch
ma V($\phi$), alors pour tout Module F , l'image inverse de F sur V($\phi$)
s'identifie, moyennant l'abus de langage habituel, au conayau de $\mathrm{id}_F\otimes\phi$
$F\otimes L^{-1} \hand{\longrightarrow} F$ ; On dit que $\phi$ est F-\uline{r\'egulier} si l'homomorphisme pr\'ec\'edent
est injectif. Si on choisit un isomorphisme de F et $\underline{O}_X$ , - ce qui
est possible localement - de sorte que $\phi$ s'identifie \`a une section de
$\underline{O}_X$ , cette terminologie estb compatible avec celle d\'ej\`a introduite par ailleur

\placedmargin{22}
